% ============================================================
%  Chapter 4 --- Indexing
% ============================================================
\chapter{Indexing}
\label{ch:indexing}

\section{Why Indexes Matter}

Without an index, answering \code{SELECT * FROM customers WHERE id = 42} requires a \emph{full table scan} --- reading every page of the table.  For a table with $N$ rows spread across $P$ pages, this is $\bigO{P}$ disk reads.  An index reduces this to $\bigO{\log P}$ (B-tree) or $\bigO{1}$ (hash index).

\begin{conceptbox}[Index = Sorted Side Structure]
An index is a \emph{separate data structure} that keeps a sorted mapping from column values to the physical locations (page ID + slot number) of the corresponding tuples in the heap file.
\end{conceptbox}

\section{Index Metadata in \shiftdb{}}

\shiftdb{} supports creating indexes via \code{CREATE INDEX} and stores them as \code{IndexMeta} records in the Catalog:

\begin{lstlisting}[style=cpp, caption={Index metadata structure.}]
struct IndexMeta {
    std::string name;
    std::string tableName;
    std::vector<std::string> columns;
    bool unique = false;
    bool isHash = false;
};
\end{lstlisting}

\begin{lstlisting}[language=SQL, caption={Creating an index.}]
CREATE INDEX idx_customers_city ON customers (city);
CREATE UNIQUE INDEX idx_email ON customers (email);
\end{lstlisting}

\section{B-Tree Fundamentals}
\label{sec:btree}

The B-tree~\cite{bayer1972btree} is the workhorse index structure of relational databases.

\subsection{Properties}

A B-tree of order $m$ satisfies:
\begin{enumerate}[nosep]
  \item Every node has at most $m$ children.
  \item Every internal node (except the root) has at least $\lceil m/2 \rceil$ children.
  \item The root has at least 2 children (unless it is a leaf).
  \item All leaves are at the same depth.
  \item A node with $k$ children contains $k-1$ keys.
\end{enumerate}

\subsection{B-Tree Structure Diagram}

\begin{figure}[H]
\centering
\begin{tikzpicture}[
  level distance=2cm,
  level 1/.style={sibling distance=6cm},
  level 2/.style={sibling distance=2.8cm},
  every node/.style={
    rectangle, draw=shiftPrimary, fill=shiftPrimary!8,
    rounded corners=2pt, font=\footnotesize,
    minimum width=2cm, minimum height=0.6cm, align=center
  }
]
\node {30 | 70}
  child { node {10 | 20}
    child { node {5 | 8} }
    child { node {12 | 15 | 18} }
    child { node {22 | 25 | 28} }
  }
  child { node {50 | 60}
    child { node {35 | 40 | 45} }
    child { node {52 | 55 | 58} }
    child { node {62 | 65 | 68} }
  }
  child { node {85 | 95}
    child { node {72 | 75 | 80} }
    child { node {88 | 90 | 92} }
    child { node {96 | 98 | 99} }
  };
\end{tikzpicture}
\caption{Example B-tree of order 4 (max 3 keys per node).}
\label{fig:btree-example}
\end{figure}

\subsection{B-Tree Operations}

\subsubsection{Search --- $\bigO{\log_m N}$}

\begin{algorithm}[H]
\caption{B-Tree Search}
\label{alg:btree-search}
\KwInput{Key $k$, root node $r$}
\KwOutput{Record pointer or \texttt{null}}
$\text{node} \leftarrow r$\;
\While{node is not a leaf}{
  Find smallest $i$ such that $k \leq \text{node.key}[i]$\;
  \If{$k = \text{node.key}[i]$}{
    \Return{\text{node.pointer}[i]}\;
  }
  $\text{node} \leftarrow \text{node.child}[i]$\;
}
Search leaf for $k$; \Return{pointer if found, else \texttt{null}}\;
\end{algorithm}

\subsubsection{Insertion}

\begin{figure}[H]
\centering
\begin{tikzpicture}[
  every node/.style={
    rectangle, draw=shiftPrimary, fill=shiftPrimary!8,
    rounded corners=2pt, font=\footnotesize,
    minimum width=2.5cm, minimum height=0.6cm, align=center
  },
  level distance=1.8cm,
  level 1/.style={sibling distance=4cm},
]

% Before split
\node[draw=none, fill=none, font=\small\bfseries] at (-5,1) {Before Insert(25):};
\node at (-5, 0) {10 | 20 | 30}
  child[draw=none] {node[draw=none, fill=none]{}} ;

% After split
\node[draw=none, fill=none, font=\small\bfseries] at (4,1) {After Insert(25):};
\node at (4,0) {20}
  child { node {10} }
  child { node {25 | 30} };
\end{tikzpicture}
\caption{B-tree node split on inserting key 25 into a full node.}
\label{fig:btree-insert-split}
\end{figure}

\begin{algorithm}[H]
\caption{B-Tree Insertion (simplified)}
\label{alg:btree-insert}
\KwInput{Key $k$, value $v$}
Find the leaf node $L$ where $k$ should be inserted\;
\If{$L$ has room}{
  Insert $(k, v)$ into $L$ in sorted order\;
}{
  Split $L$ into $L'$ and $L''$ around the median key\;
  Push median key up to parent\;
  \If{parent overflows}{
    Recursively split upward\;
  }
}
\end{algorithm}

\subsubsection{Deletion}

\begin{algorithm}[H]
\caption{B-Tree Deletion (simplified)}
\label{alg:btree-delete}
\KwInput{Key $k$}
Find the leaf node $L$ containing $k$\;
Remove $k$ from $L$\;
\If{$L$ has fewer than $\lceil m/2 \rceil - 1$ keys}{
  \If{sibling can donate a key}{
    Redistribute (borrow from sibling through parent)\;
  }{
    Merge $L$ with sibling, pulling separator from parent\;
    \If{parent underflows}{
      Recursively fix upward\;
    }
  }
}
\end{algorithm}

\section{B$^+$-Tree vs B-Tree}

\begin{table}[H]
\centering
\caption{B-tree vs B$^+$-tree comparison.}
\label{tab:btree-vs-bplus}
\begin{tabular}{l p{6cm} p{6cm}}
\toprule
\textbf{Property} & \textbf{B-Tree} & \textbf{B$^+$-Tree} \\
\midrule
Data pointers      & In all nodes      & Only in leaves \\
Leaf linkage       & None              & Leaves linked (range scans) \\
Internal node size & Larger (has data) & Smaller (more keys per node) \\
Range queries      & Slow (traverse tree) & Fast (leaf chain scan) \\
Used by            & Textbook examples  & PostgreSQL, MySQL/InnoDB, SQLite \\
\bottomrule
\end{tabular}
\end{table}

\section{Hash Indexes}
\label{sec:hash-index}

For \emph{equality-only} lookups, a hash index provides $\bigO{1}$ amortised access.

\begin{conceptbox}[Static vs Dynamic Hashing]
\begin{itemize}[nosep]
  \item \textbf{Static hashing}: fixed number of buckets --- simple but suffers from overflow chains.
  \item \textbf{Extendible hashing}: uses a directory that doubles in size on overflow --- no overflow chains but more complex.
  \item \textbf{Linear hashing}: splits buckets one at a time in round-robin order --- no directory needed.
\end{itemize}
\end{conceptbox}

\section{Index Selection Guidelines}

\begin{table}[H]
\centering
\caption{When to use which index type.}
\begin{tabular}{l c c l}
\toprule
\textbf{Workload} & \textbf{B-Tree} & \textbf{Hash} & \textbf{Reason} \\
\midrule
Point queries (\code{=})              & \cmark & \cmark & Both work; hash is faster \\
Range queries (\code{<}, \code{BETWEEN}) & \cmark & \xmark & Hash cannot range scan \\
\code{ORDER BY}                        & \cmark & \xmark & B-tree leaves are sorted \\
\code{LIKE 'abc\%'}                    & \cmark & \xmark & Prefix matches use tree \\
Frequent updates                       & \cmark & \cmark & Both handle updates \\
\bottomrule
\end{tabular}
\end{table}

\section{Complexity Summary}

\begin{table}[H]
\centering
\caption{Asymptotic costs of index operations (order-$m$ B-tree, $N$ keys).}
\begin{tabular}{l c c c}
\toprule
\textbf{Operation} & \textbf{Heap Scan} & \textbf{B-Tree} & \textbf{Hash} \\
\midrule
Search  & $\bigO{N}$ & $\bigO{\log_m N}$ & $\bigO{1}$ amortised \\
Insert  & $\bigO{1}$ & $\bigO{\log_m N}$ & $\bigO{1}$ amortised \\
Delete  & $\bigO{N}$ & $\bigO{\log_m N}$ & $\bigO{1}$ amortised \\
Range   & $\bigO{N}$ & $\bigO{\log_m N + k}$ & $\bigO{N}$ \\
\bottomrule
\end{tabular}
\end{table}

\begin{interviewbox}[Interview Corner]
\textbf{Q:} Why do most production databases use B$^+$-trees instead of B-trees?\\[4pt]
\textbf{A:} B$^+$-trees store data only in leaf nodes, so internal nodes are smaller and more keys fit per page --- this reduces the tree height.  Leaves are linked in a doubly-linked list, making range scans sequential rather than requiring tree traversals.

\vspace{6pt}
\textbf{Q:} What is a \emph{clustered} index?\\[4pt]
\textbf{A:} A clustered index determines the \emph{physical order} of rows on disk.  A table can have at most one clustered index.  In MySQL/InnoDB, the primary key is always the clustered index (called IOT --- index-organised table).

\vspace{6pt}
\textbf{Q:} When would you choose a hash index over a B-tree?\\[4pt]
\textbf{A:} When the workload is dominated by exact-match lookups (\code{WHERE id = 42}) with no range queries, no \code{ORDER BY}, and no prefix \code{LIKE} searches.  Hash indexes provide $\bigO{1}$ lookups but cannot serve ranges.

\vspace{6pt}
\textbf{Q:} What is the cost of inserting into a B-tree?  What triggers a split?\\[4pt]
\textbf{A:} Insertion is $\bigO{\log_m N}$ because you traverse the tree to find the correct leaf, then insert.  A split occurs when a node exceeds $m-1$ keys; the median key is pushed to the parent.  If the parent also overflows, the split cascades upward --- in the worst case, the root splits and the tree height increases by one.
\end{interviewbox}
